\subsection{Vlerësuesi, animi dhe gabimi kuadratik mesatar}
Një vlerësues $\widehat\theta$ është ndryshore rasti sepse varet nga mostra. Animi është
\begin{equation}
 \operatorname{Bias}(\widehat\theta)=\E[\widehat\theta]-\theta.
\end{equation}
Gabimi kuadratik mesatar ndahet si
\begin{equation}
 \E[(\widehat\theta-\theta)^2]
 =\Var(\widehat\theta)+\operatorname{Bias}(\widehat\theta)^2.
\end{equation}
Ky identitet tregon se vlerësuesi i paanshëm nuk është automatikisht optimal: një anim i vogël mund të pranohet nëse ul ndjeshëm variancën. Në simulimet fizike, kjo shfaqet te rregullimi i modeleve, prerja e peshave ekstreme ose vlerësuesit me kontroll variabël.

Varianca e mostrës me emërues $N-1$ është e paanshme për variancën e popullatës kur mostrat janë të pavarura,
\begin{equation}
 s^2=\frac{1}{N-1}\sum_{i=1}^{N}(X_i-\overline X)^2.
\end{equation}
Emëruesi $N-1$ vjen sepse mbetjet kanë shumën zero dhe vetëm $N-1$ shkallë lirie. Për seri të korreluara, formula ende përshkruan shpërndarjen margjinale, por $s/\sqrt N$ nuk është gabimi standard i mesatares.

\subsubsection{Intervalet e besimit dhe interpretimi}
Një interval besimi $95\%$ nuk do të thotë se, pasi të dhënat janë vëzhguar, parametri fiks ka probabilitet $95\%$ të jetë brenda atij intervali. Interpretimi frekuentist është se procedura prodhon intervale që mbulojnë parametrin në $95\%$ të përsëritjeve ideale. Ky dallim është i rëndësishëm në raportim.

Për një mesatare me variancë të panjohur dhe të dhëna Gauss të pavarura,
\begin{equation}
 \frac{\overline X-\mu}{s/\sqrt N}\sim t_{N-1}.
\end{equation}
Për mostra të mëdha, kuantilet $t$ afrohen me ato normale. Për probabilitete binomiale, intervali Wilson ose një interval i bazuar në gjasë është më i qëndrueshëm se intervali Wald, sidomos pranë kufijve.

\subsubsection{Bootstrap-i}
Bootstrap-i joparametrik trajton shpërndarjen empirike si përafrim të popullatës. Nga mostra origjinale me $N$ elemente krijohen mostra të reja me rizgjedhje dhe të njëjtën madhësi. Për secilën llogaritet statistiku $T^{*(b)}$. Shpërndarja e këtyre statistikëve vlerëson variancën, animin dhe intervalet.

Metoda funksionon mirë për shumë statistikë të lëmuar, por nuk është universale. Për maksimumin, bishtat ekstremë dhe mostra shumë të vogla mund të jetë e dobët. Për seri kohore nuk duhet rizgjedhur çdo pikë veçmas, sepse prishet varësia; përdoret bootstrap në blloqe. Në Monte Carlo me peshë, skema e rizgjedhjes duhet të respektojë mënyrën si u krijuan peshat.

\subsubsection{Kovarianca nga burime të përbashkëta}
Në eksperimente fizike, dy madhësi shpesh varen nga i njëjti kalibrim. Nëse $x=x_0(1+\epsilon)$ dhe $y=y_0(1+\epsilon)$ ndajnë të njëjtin gabim relativ, raporti $x/y$ anulon gabimin e përbashkët në rendin e parë. Trajtimi i $x$ dhe $y$ si të pavarura do ta mbivlerësonte pacaktueshmërinë. Në diferencën $x-y$, i njëjti korrelacion mund ta ulë ose ta rrisë efektin në varësi të shenjave.

Matrica e kovariancës duhet të jetë simetrike pozitivisht gjysmë e përcaktuar. Një matricë e ndërtuar nga koeficientë korrelacioni të papajtueshëm mund të ketë vlera vetjake negative dhe të mos përfaqësojë asnjë shpërndarje. Para kampionimit shumënormal kontrollohen simetria dhe spektri; vlera të vogla negative nga rrumbullakimi mund të kërkojnë korrigjim të dokumentuar.

\subsubsection{Ligji total i variancës}
Kur modeli varet nga parametri i rastit $\Theta$ dhe nga zhurma e brendshme,
\begin{equation}
 \Var(Y)=\E[\Var(Y\mid\Theta)]
 +\Var(\E[Y\mid\Theta]).
\end{equation}
Termi i parë mat variabilitetin brenda realizimeve për parametra të fiksuar; i dyti mat ndryshimin ndërmjet parametrave. Kjo ndarje është e dobishme në simulimin e zbritjes së mjetit hapësinor, ku mund të dallohen ndryshimet atmosferike nga zhurma e sensorit ose e kontrollit.
